\section{Linear Search: Completion, Implementation, and Evidence}\label{ec:r36-linear}
This section completes the implementation details of Theorem~\ref{thm:r36-linear}. All economic reductions, the continuous terminal minimization, and the moment arrays remain those in Sections~\ref{sec:r34-global} and~\ref{sec:r35-monge}. The additional step is an ordinal matrix interface on which every submatrix has ordered row minima. It is not enough to verify only the minima of the original triangular matrix.

\subsection{Nested feasible domains and exact ties}
For each fixed search, use strictly ordered row and column indices. A terminal row $i$ has feasible columns $\ell\geq i$. A prefix row with endpoint $j$ has feasible predecessor indices $i<j$, after columns that cannot be reached at the preceding layer have been removed. The first prefix layer has only its single initialized predecessor. Subsequent layers restrict the columns to the appropriate suffix of globally reachable predecessor indices. Every production row then has a feasible column.

For a feasible entry, let $K_{rc}=(0,A_{rc},c)$. On an upper staircase, set $K_{rc}=(1,0,-c)$ for a missing prefix entry. On a lower staircase, set $K_{rc}=(1,0,c)$ for a missing suffix entry. Compare keys lexicographically. Column indices make all keys within a row distinct, including equal-cost feasible entries. Feasible entries always precede missing entries. The completion is ordinal: it assigns neither a fabricated rational penalty nor a floating-point infinity to an unavailable contract.

For clarity, we verify the pairwise single-crossing condition in all cases. Fix rows $r<s$ and columns $c<d$ and suppose $K_{rd}<K_{rc}$. When both compared columns are feasible in both rows, this preference means $A_{rd}<A_{rc}$. The feasible Monge inequality yields
\[
 A_{sd}-A_{sc}\leq A_{rd}-A_{rc}<0,
\]
so the preference persists. The same calculation applies after adding an arbitrary preceding-layer potential to each column. Transposition of the anchored-edge or deterministic-cell array preserves the relevant Monge inequality for endpoint rows and predecessor columns.

In the upper staircase there are three other possibilities. If $c$ is missing and $d$ is feasible, the earlier preference follows from the first component. In a later row either this status persists or both entries become missing, in which case $-d<-c$. If both earlier entries are missing, both later entries are missing, with the same strict preference. If both earlier entries are feasible but one or both later entries become missing, those same two cases apply. A feasible $c$ and missing $d$ cannot occur on this staircase.

On the lower staircase a missing $d$ cannot be preferred to feasible $c$, and when both are missing the positive column indices prefer $c$. Thus an earlier preference for $d$ requires both entries to be feasible. Since feasible prefixes expand with row index, both remain feasible in the later row and the finite Monge argument applies. This exhausts the cases. The implication survives deletion of any rows or columns, including deletion of every feasible entry of an original row. It proves total monotonicity with a unique ordinal minimizer in every nonempty submatrix. Only the full production rows must have feasible minima.

A two-column regression explains why the signs matter. Restrict an upper staircase to columns 1 and 2 and rows 2 and 3 with zero finite costs. At row 2 only column 2 is feasible. At row 3 neither is feasible. Uniform missing values with leftmost ties yield minimizing columns 2 and 1, violating monotonicity. The negative-index prefix keys instead yield columns 2 and 2. This is a diagnostic of the linear-search interface, not a counterexample to the feasible Monge theorem or to the earlier divide-and-conquer algorithm.

\subsection{Reduction and interpolation invariants}
The function \texttt{smawk\_minima} implements the classical reduction, recursion, and interpolation procedure. For an ordered row list $R$ and column list $C$, scan $C$ from left to right while maintaining a stack of at most $|R|$ columns. Compare an incoming column with the stack's last column at the row whose position is the current stack length minus one. Pop the last column precisely when the incoming key is smaller. Append the incoming column if the stack is shorter than the row list.

The reduction preserves a minimizing column for every retained row. Suppose a stack column at zero-based position $t$ is popped by a new column. At witness row $R_t$ the new column is preferred, hence it dominates the popped column at all later rows by single crossing. If $t>0$, the preceding stack column was preferred to the popped column at $R_{t-1}$ when the latter was appended. By the contrapositive of single crossing it dominates at every earlier row. That predecessor cannot have changed while its successor remained on the stack. Thus the popped column is never uniquely needed; later pops only replace its dominators. At $t=0$ there are no earlier rows. If the stack is full and an incoming column is not appended, the last stack column beats it at the final witness and consequently at every earlier row. This proves reduction, including missing-only submatrices, without a numerical assumption about the padding. It is the classical matrix reduction, not a claim inferred from tests.

Recursively solve the odd-positioned rows using the reduced column list. For each remaining even-positioned row, scan only between the minimizing columns of its solved neighbors, using the full reduced-list boundary where a neighbor is absent. Total monotonicity places its minimizer in that range. A position map is constructed once per recursion level; the code does not repeatedly perform a linear search for a column's position. Every comparison uses the complete exact key, including tie-breaking indices.

Column reduction uses $O(|R|+|C|)$ comparisons. Interpolation scans ranges whose interiors do not overlap, apart from neighboring boundaries, and hence costs $O(|R|+|C'|)$, where $C'$ is the reduced list. Recursion halves the row count, and after the first reduction the column count is no larger than the preceding row count. Summing over levels gives $O(|R|+|C|)$ comparisons and stored indices. No cache of all feasible matrix entries is constructed. Winner values and terminal levels are recomputed after search; this adds only linear work and is counted in the reported oracle totals.

The implementation uses an explicit key callback and ordered index lists. Its precondition is total monotonicity; it does not attempt an expensive automatic proof for arbitrary supplied cost functions. Invalid or repeated row/column indices are rejected. Empty row sets return an empty result; a nonempty row set with no columns is rejected. The staircase wrapper rejects a missing final minimum, which signals a violated feasibility precondition rather than silently converting a missing entry into a contractual value.

\subsection{Arithmetic, storage, and execution accounting}
The six prefix-moment arrays require linear preprocessing for already ordered inputs. Each finite terminal query constructs its exact quadratic, clamps the rational stationary point to the interval, and returns its objective and level. The singleton last interval is evaluated directly by the same interface. Each finite prefix or deterministic cell uses constant-many reduced rational operations. Lexicographic padding and index comparisons require $O(\log(k+1))$ bits, which are absorbed in the coefficient-height bound $S$.

There are one terminal search and at most $m$ program-layer searches. Combining terminal and preceding-prefix costs requires a linear scan per budget; only the winning codebook is reconstructed. All codebooks together have at most $1+\cdots+m$ levels, so reconstruction and output storage are $O(m^2)\subseteq O(mk)$. Backpointers use $O(mk)$ indices, while the active search, two cost rows, and terminal optima use $O(k)$ space. As in the preceding arithmetic model, indices and pointers occupy $O(\log(k+1))$ bits. The conservative Turing bound charges comparisons, integer products, and reduced-rational arithmetic; it does not equate a recorded oracle call with one bit operation. Unordered input additionally requires sorting, with $O(k\log(k+1))$ rational comparisons. Budgets above $k$ append references to the attained zero-loss solution; the theorem is stated for $m\leq k$ so that output length is explicitly charged.

The output per budget is a codebook. In the expected-participation institution, the terminal service draw is an adjacent-codeword lottery and the intermediate service makes the branch's conditional-mean payment equal its cap. Fresh randomness is not retained as an uncharged state. In the pathwise institution, the same branch uses its largest feasible deterministic level; Theorem~\ref{thm:r34-pathwise} identifies this as the unrestricted randomized pathwise optimum in the stated renewal architecture. Search does not change either observation interface or admissibility definition.

\subsection{Independent checks and comparator scope}
The R36 suite uses 100 exact rational instances, including one branch, equalities in marginal monotonicity, zero intermediate curvature, highly unequal probabilities, and budgets exceeding the exact alphabet threshold. It records 4,232 objective equalities against the unchanged R34 and R35 algorithms and 2,116 controller replays. An independent comparator enumerates every theorem-permitted anchored prefix and obtains each last-interval quadratic from three direct branchwise evaluations, rather than the production moments or dynamic program. It checks 156 budget objectives across 550 terminal candidates. These are exhaustive checks of the theorem-reduced candidates, not an enumeration of the continuum.

The search-specific checks examine all nonempty row and column subsets for selected small actual and tied staircase matrices, including rows whose feasible columns have all been deleted. There are 43,431 exact row-argmin comparisons, 2,595 nonvacuous single-crossing implications, and 16,273 conservative work-bound checks. Additional checks cover 255 directly reconstructed terminal cells, 402 feasible terminal Monge inequalities, 176 nonanchor grid candidates, two explicit padding regressions, and 17 rejected invalid inputs. Nonanchor grids are falsification tests, not the basis of global optimality.

The same-input study computes all eight requested budgets at each size from 16 through 2,048 branches. Both frontiers agree with the preceding monotone implementation at every size; the quadratic randomized algorithm is rerun through 256 branches. Larger quadratic cases are unexecuted, not censored timing results. Complete codebooks, exact losses, work counts, actual runtimes, environment metadata, and comparator source hashes are saved. Repeated oracle evaluations count repeatedly. The ordinal comparisons and economic queries are reported separately.

The linear bound is an asymptotic guarantee, not uniform practical dominance. At 16 branches, the new expected-participation search makes 534 economic queries versus 289 for divide-and-conquer; at 2,048 branches it makes 128,966 versus 135,597. The deterministic counts at the latter size are 163,838 versus 153,180, so its crossover has not occurred in this study. These observations are retained, not filtered out. The three versions solve the identical global contractual frontier, but share the economic reduction and moment interface. No comparison with a published external parametric-flow solver or calibrated operational dataset is claimed.
